\documentclass[11pt,a4paper]{article}
\usepackage[utf8]{inputenc}
\usepackage{amsmath}
\usepackage{amsfonts}
\usepackage{amssymb}
\usepackage{graphicx}
\usepackage{hyperref}
\usepackage{geometry}
\usepackage{booktabs}
\usepackage{xcolor}
\usepackage{enumitem}
\usepackage{fancyhdr}
\geometry{a4paper, margin=0.9in}

\hypersetup{colorlinks=true, linkcolor=blue, urlcolor=blue, citecolor=blue}

\definecolor{tier1}{HTML}{2E86AB}
\definecolor{tier2}{HTML}{A23B72}
\definecolor{tier3}{HTML}{F18F01}
\definecolor{tier4}{HTML}{C73E1D}

\title{\textbf{Research Roadmap: opto-sim Physical-Layer QKD Simulator}\\[6pt]
\large A Tiered Plan with Effort Estimates and Diminishing-Returns Analysis}
\author{Based on the opto-sim codebase --- \texttt{D:\textbackslash Azaan\textbackslash Study\textbackslash git\textbackslash opto-sim-dev}}
\date{\today}

\begin{document}

\maketitle

% ---- OUTLINE PAGE ----
\subsection*{What makes this simulator different}
Most QKD simulators commit one of two sins.\\
\\
\textbf{Sin 1 --- Abstract qubit layer}: NetSquid, SimulaQron, and the QuTIP-based simulators represent quantum states as vectors
$|\psi\rangle$ and channels as depolarising maps.  A ``detector'' clicks with some probability.  There is no field, no wavelength,
no phase noise, no birefringence, no \emph{physics} between Alice's laser and Bob's APD.\\
\\
\textbf{Sin 2 --- Numeric microwave}: Lumerical, RSoft, and OptiSystem are exquisite for photonic circuit design but are
closed-source, priced per seat, and designed for sinusoidal carriers---not for QKD, where the complex envelope carries the
quantum information.\\
\\
\textbf{opto-sim is neither}.  It propagates the complex-envelope electric field (power-bearing, calibrated once at the laser)
through physically parametrised components: a CWLaser with Wiener phase noise and relaxation-oscillation RIN, a PhaseModulator
whose $V_\pi$ comes from LiNbO$_3$ crystal geometry, a MZM built from two modulators in an interferometer, a fibre whose
birefringence \emph{depends on temperature and bend radius} via analytic formulae, and an APD whose shot noise carries the
excess-noise factor from the gain process.

\vspace{6pt}
\noindent This permits research questions that the other simulators simply cannot address:
\begin{itemize}[nosep]
    \item How does a \$10 temperature swing on an aerial fibre change the sifted key rate?
    \item Can the relaxation-oscillation peak of the RIN spectrum leak information to an eavesdropper?
    \item What is the joint effect of CD + PMD + polarisation drift on BB84 phase encoding at 10 Gbaud?
\end{itemize}

\newpage
% ---- THE PLAN ----
\section{The Tiered Research Plan}

Each tier is self-contained and publishable.  The effort column is for \emph{one} researcher
working in this existing codebase.

\bigskip

% --- TIER 0 ---
\subsection*{\color{tier1}Tier 0 \hfill Reproducibility \& testing infrastructure}
\addcontentsline{toc}{subsection}{Tier 0: Reproducibility \& testing infrastructure}

\begin{tabular}{ll}
\toprule
\textbf{Effort} & 1--2 weeks \\
\textbf{Deliverable} & pytest suite + CI + seeded-reproducibility \\
\textbf{Paper type} & None (scaffolding for all tiers) \\
\bottomrule
\end{tabular}

\vspace{4pt}
\noindent\textbf{What to do:}
\begin{enumerate}[nosep]
    \item Pin a fixed RNG seed per test so every ``random'' fibre realisation is repeatable.
    \item Write unit tests for each component (CWLaser power~\& linewidth, MZM transfer curve, CD pulse broadening, APD noise variance).
    \item Add a \texttt{tests/} directory with \texttt{conftest.py} and a \texttt{Makefile} / \texttt{tox.ini}.
    \item Add a \texttt{--seed} CLI argument to BB84 scripts.
\end{enumerate}
\noindent\textbf{Why first:} Without seeded randomness, every figure in every paper will be unreproducible.
Every hour spent here saves days of debugging false QBER changes later.

\vspace{4pt}
\noindent\textbf{Diminishing returns:} Stop after you can regenerate \texttt{qber\_vs\_distance.png}
bit-for-bit identical from \texttt{git checkout}.  Do not build a full CD/CI matrix yet.

\newpage

% --- TIER 1 ---
\subsection*{\color{tier1}Tier 1 \hfill Channel validation and characterisation}
\addcontentsline{toc}{subsection}{Tier 1: Channel validation and characterisation}

\begin{tabular}{ll}
\toprule
\textbf{Effort} & 2--4 weeks \\
\textbf{Deliverable} & Validated fibre model against published data \\
\textbf{Paper type} & Conference (OFC, CLEO, ECOC) \\
\bottomrule
\end{tabular}

\vspace{4pt}
\noindent\textbf{What makes this novel:} No QKD simulator validates chromatic dispersion, PMD, and
birefringence against separate experimental datasets.  Doing so makes the simulator defensible
in peer review.

\vspace{4pt}
\noindent\textbf{Tasks:}

\begin{enumerate}[nosep]
    \item \textbf{CD validation} (3 d).  Reproduce Agrawal Fig 2.6 (Gaussian pulse broadening at $z/L_D = 0.5, 1.0, 2.0$).
          Fit $\beta_2$ from SMF-28 datasheet (D = 17 ps/(nm\,km)) and confirm width ratio within 0.1\,\%.

    \item \textbf{PMD validation} (3 d).  Generate 10$^4$ fibre realisations, plot DGD histogram against
          Maxwellian PDF (Razavi [5] Fig 2.11).  Compute $\langle \Delta\tau \rangle$ and confirm it matches
          the PMD coefficient $\times \sqrt{L}$.

    \item \textbf{Attenuation validation} (1 d).  Plot received power vs distance for an SMF-28 link
          at 1550 nm (0.182 dB/km).  Overlay OTDR data from Keiser [1] or a Thorlabs application note.

    \item \textbf{Birefringence validation} (2 d).  Reproduce Yuan [4] Fig 1: beat length $L_B$ vs fibre
          bend radius for various $\Delta n_0$.  Confirm $L_B \propto R / \Delta n_0$ scaling.

    \item \textbf{Reproducibility pass} (2 d).  All scripts save PNG~\& CSV with \texttt{(seed,params)} in the
          filename.  A reader can re-run with \texttt{python run\_all.py --seed 42}.
\end{enumerate}

\vspace{4pt}
\noindent\textbf{Effort breakdown:} 11 d coding, 3 d writing, 1 d figures = \textbf{3 weeks}.

\vspace{4pt}
\noindent\textbf{Diminishing returns:} Once each model reproduces its reference figure within
reasonable tolerance, move on.  Spending weeks to shave the last 0.01\,\% off the CD fit
yields zero research value --- the question is whether the model is \emph{physically correct},
not whether it is perfect.

\newpage

% --- TIER 2 ---
\subsection*{\color{tier2}Tier 2 \hfill Environmental sensitivity of QKD}
\addcontentsline{toc}{subsection}{Tier 2: Environmental sensitivity of QKD}

\begin{tabular}{ll}
\toprule
\textbf{Effort} & 4--8 weeks \\
\textbf{Deliverable} & Temperature/bend $\rightarrow$ QBER curves, polarisation drift analysis \\
\textbf{Paper type} & Journal (Optics Express, IEEE JLT, EPJ QT) \\
\bottomrule
\end{tabular}

\vspace{4pt}
\noindent\textbf{What makes this novel:} The birefringence formula in \texttt{fiber.py}
(Yuan [4]) makes temperature and bend radius \emph{explicit} parameters.  Every other QKD
simulator treats polarisation as a random unitary.  This model can answer: \emph{would burying
the fibre vs aerial deployment change the QBER statistics measurably?}

\vspace{4pt}
\noindent\textbf{Tasks:}

\begin{enumerate}[nosep]
    \item \textbf{Temperature sweep} (1 w).  0--50 $^\circ$C, 25 km spans.  Plot $\Delta n(T)$,
          beat length, and resulting QBER for BB84 at 1 Gbaud.  Does the phase shift $\theta(T)$
          cross a zero (all bits flip)?  How fast?

    \item \textbf{Bend sensitivity} (1 w).  \(R = 10\)--100 mm, fixed $T$.  Plot $\Delta n(R)$
          via Yuan Eq 8--9.  Coupled with temperature: at what (T, R) combination does the
          fibre become visibly polarisation-maintaining ($L_B < 1$ m)?

    \item \textbf{Time-varying channel} (2 w).  Inject a slow temperature ramp
          ($1^\circ$C/min) and run repeated BB84 exchanges.  Plot QBER as a time series.
          Does the reconciliation protocol recover?  How large a temperature swing is tolerable
          before the key rate drops below zero?

    \item \textbf{Aerial vs buried comparison} (1 w).  Use published weather data (e.g., ASHRAE
          handbook) to bound diurnal temperature variation.  Compare QBER variance for aerial
          (fast, large swings) vs buried (slow, damped) scenarios.

    \item \textbf{Write-up} (1 w).  Figures, discussion, comparison to literature.
\end{enumerate}

\vspace{4pt}
\noindent\textbf{Effort breakdown:} 5 w coding + simulation, 1 w writing, 1 w review = \textbf{7 weeks}.

\vspace{4pt}
\noindent\textbf{Diminishing returns:} The environmental story is strongest when it makes a
concrete, falsifiable prediction (e.g., ``QBER variance in aerial fibre is 3$\times$ larger
than buried'').  Do not simulate 50 scenarios; find the threshold and stop.

\newpage

% --- TIER 3 ---
\subsection*{\color{tier3}Tier 3 \hfill Physical-layer security analysis}
\addcontentsline{toc}{subsection}{Tier 3: Physical-layer security analysis}

\begin{tabular}{ll}
\toprule
\textbf{Effort} & 3--5 months \\
\textbf{Deliverable} & Side-channel quantification, finite-key analysis \\
\textbf{Paper type} & High-impact journal (IEEE T-IFS, NPJ Quantum Inf, PRX Quantum) \\
\bottomrule
\end{tabular}

\vspace{4pt}
\noindent\textbf{Why it matters:} The physical-layer approach lets you simulate attacks that
abstract simulators cannot even represent.  This is where the simulator earns its keep.

\vspace{4pt}
\noindent\textbf{Tasks:}

\begin{enumerate}[nosep]
    \item \textbf{Laser noise side-channel} (3 w).  CWLaser's RIN resonance (Coldren Eq 5.3.38)
          produces correlated amplitude fluctuations.  Can Eve synchronise to the relaxation peak
          and extract bit information from the amplitude envelope?  Quantify leaked information
          in bits/symbol.

    \item \textbf{Finite-key effects with realistic noise} (4 w).  Derive SKR from simulated
          QBER using the canonical GLLP + decoy-state formulas.  Show how finite-key bounds
          change when the underlying noise is correlated (RIN, 1/f phase noise) rather than
          i.i.d.  This is a genuinely open question in QKD security proofs.

    \item \textbf{Detector side-channel} (2 w).  APD dead time and afterpulsing probabilities
          create detector-efficiency mismatches.  Simulate time-shift attacks (Makarov 2009)
          in the physical-layer framework.  Confirm the attack success rate vs QBER penalty.

    \item \textbf{Phase modulator flaw analysis} (2 w).  MZM drive voltage ripple
          ($V = V_\pi + \delta\sin(\omega t)$) creates a deterministic phase error pattern.
          Can Eve exploit this pattern to learn Alice's basis choice?

    \item \textbf{Write-up} (3 w).  A full journal paper.
\end{enumerate}

\vspace{4pt}
\noindent\textbf{Effort breakdown:} 3 m coding, 1 m writing, 1 m review = \textbf{5 months}.

\vspace{4pt}
\noindent\textbf{Diminishing returns:} Do not implement every attack in the literature.
Pick the one where the physical-layer model gives a clear advantage (the laser-noise
side-channel or the correlated-noise finite-key analysis).  One strong result beats a
catalogue of re-runs.

\newpage

% --- TIER 4 ---
\subsection*{\color{tier4}Tier 4 \hfill Full QKD system optimisation}
\addcontentsline{toc}{subsection}{Tier 4: Full QKD system optimisation}

\begin{tabular}{ll}
\toprule
\textbf{Effort} & 6--12 months \\
\textbf{Deliverable} & Decoy-state protocol, SKR optimisation, open-source release \\
\textbf{Paper type} & Nature Communications, PRX Quantum + code release \\
\bottomrule
\end{tabular}

\vspace{4pt}
\noindent\textbf{Why it matters:} This closes the loop: physical model $\rightarrow$ protocol
$\rightarrow$ secure key rate $\rightarrow$ optimisation.  A well-documented open-source
framework with a reproducible paper attached would be unique in the QKD simulation landscape.

\vspace{4pt}
\noindent\textbf{Tasks:}

\begin{enumerate}[nosep]
    \item \textbf{Decoy-state protocol} (2 m).  Implement 3-intensity decoy-state BB84
          (Lo, Ma \& Chen 2005).  Optimise $\mu, \nu, \omega$ intensities for the physical
          laser model.  Show how the finite-intensity statistical fluctuations interact with
          actual RIN to change the secure bound.

    \item \textbf{SKR optimisation} (2 m).  Multi-parameter sweep: laser power, modulation
          depth, detector bandwidth, fibre length, temperature.  Produce contour plots of
          SKR in (distance, power) space.  Include detector dead-time limitation at high
          rates.

    \item \textbf{Framework clean-up and documentation} (1 m).  API docs, tutorial notebooks,
          a one-command demo.  Publish to PyPI or as a GitHub release with a DOI.

    \item \textbf{Write-up} (2 m).  The definitive paper showing what a physical-layer QKD
          simulator can do.
\end{enumerate}

\vspace{4pt}
\noindent\textbf{Effort breakdown:} 6 m coding, 2 m writing, 2 m review = \textbf{10 months}.

\vspace{4pt}
\noindent\textbf{Diminishing returns:} This tier is where feature creep is most dangerous.
Resist the urge to add every protocol (E91, MDI-QKD, TF-QKD), every detector (SNSPD, TES),
and every modulation format (PPM, DPSK, CSK).  The goal is \emph{one} complete system story,
not a tool chest.  Adding more protocols past the first is a \textbf{strongly diminishing}
activity.

\newpage

\section{Summary and Recommended Route}

\begin{center}
\begin{tabular}{lccr}
\toprule
\textbf{Tier} & \textbf{Time} & \textbf{Impact} & \textbf{Risk} \\
\midrule
0: Reproducibility    & 1--2 w   & Scaffolding    & None \\
1: Channel validation & 3 w      & Conference     & Low \\
2: Environmental      & 7 w      & Journal        & Medium \\
3: Security           & 5 m      & High-impact J. & Medium-High \\
4: Full system        & 10 m     & Flagship       & High \\
\bottomrule
\end{tabular}
\end{center}

\vspace{12pt}
\noindent\textbf{Recommended strategy:}
\begin{enumerate}[nosep]
    \item Do Tier 0 first (2 w).  Every tier needs it.
    \item Publish Tier 1 at OFC/CLEO (3 w).  This is a low-risk calibration paper.
    \item If Tier 2 yields a clean result, publish at Optics Express (7 w).
          If the environmental story is weak, skip to Tier 3.
    \item Tier 3 is the highest impact per hour spent for this codebase ---
          \emph{no other simulator can do the laser-noise side-channel analysis}.
          Invest 5 months here.
    \item Tier 4 only if the project needs a flagship capstone.  The marginal impact
          of the tenth contour plot is much smaller than the first novel attack.
\end{enumerate}

\vspace{12pt}
\noindent\textbf{The one-page version:}

\vspace{6pt}
\noindent\fcolorbox{gray}{white}{\parbox{0.95\textwidth}{
\textbf{Do this first:} Tier 0 + Tier 1 (5 weeks total).\\
\textbf{Then publish} a channel-validation conference paper.\\
\textbf{Then do} the laser-noise side-channel from Tier 3 (3 months).\\
\textbf{Stop there.} You will have two papers and a defensible, validated simulator ---
the diminishing returns past this point are steep.
}}

\end{document}
